%%
%% This is file `sample-sigconf.tex',
%% generated with the docstrip utility.
%%
%% The original source files were:
%%
%% samples.dtx  (with options: `sigconf')
%% 
%% IMPORTANT NOTICE:
%% 
%% For the copyright see the source file.
%% 
%% Any modified versions of this file must be renamed
%% with new filenames distinct from sample-sigconf.tex.
%% 
%% For distribution of the original source see the terms
%% for copying and modification in the file samples.dtx.
%% 
%% This generated file may be distributed as long as the
%% original source files, as listed above, are part of the
%% same distribution. (The sources need not necessarily be
%% in the same archive or directory.)
%%
%% The first command in your LaTeX source must be the \documentclass command.
\documentclass[sigconf]{acmart}

\usepackage[T1]{fontenc}
% \usepackage[polish]{babel}
\usepackage[utf8x]{inputenc}

\usepackage{comment}
\usepackage{minted}
\usepackage{xcolor}
\usepackage{multirow}
\usepackage{breqn}
\usepackage{mathptmx} % This is Times font
\usepackage{fancyhdr}
\usepackage{mathrsfs}
\usepackage{multirow}
\usepackage{listings}
\usepackage{graphicx}
\usepackage[export]{adjustbox}
\usepackage{caption}
\usepackage{subcaption}
\usepackage{setspace}
\usepackage{todonotes}
\usepackage{array}
\usepackage{multirow}
\usepackage{bm}
% \usepackage[linesnumbered, ruled, lined, boxed]{algorithm2e}
% \usepackage[ruled, lined, boxed]{algorithm2e}
\usepackage[ruled,vlined,linesnumbered]{algorithm2e}
\usepackage{algpseudocode}
\usepackage{booktabs}
\usepackage{threeparttable}
\usepackage{enumerate}
\usepackage{color}
\usepackage{xcolor,colortbl}
\usepackage{adjustbox}
\usepackage{soul}
\usepackage[inline]{enumitem}
\usepackage{pifont}
\usepackage{float}
% \renewcommand\thesubfigure{(\alph{subfigure})}

\floatstyle{plaintop}
\restylefloat{table}
\SetArgSty{textnormal}
% \setcounter{secnumdepth}{6}

\DeclareMathOperator*{\argmax}{arg\,max}
\DeclareMathOperator*{\argmin}{arg\,min}


%% NOTE that a single column version may be required for 
%% submission and peer review. This can be done by changing
%% the \doucmentclass[...]{acmart} in this template to 
%% \documentclass[manuscript,screen]{acmart}
%% 
%% To ensure 100% compatibility, please check the white list of
%% approved LaTeX packages to be used with the Master Article Template at
%% https://www.acm.org/publications/taps/whitelist-of-latex-packages 
%% before creating your document. The white list page provides 
%% information on how to submit additional LaTeX packages for 
%% review and adoption.
%% Fonts used in the template cannot be substituted; margin 
%% adjustments are not allowed.
%%
%%
%% \BibTeX command to typeset BibTeX logo in the docs
\AtBeginDocument{%
  \providecommand\BibTeX{{%
    \normalfont B\kern-0.5em{\scshape i\kern-0.25em b}\kern-0.8em\TeX}}}

%% Rights management information.  This information is sent to you
%% when you complete the rights form.  These commands have SAMPLE
%% values in them; it is your responsibility as an author to replace
%% the commands and values with those provided to you when you
%% complete the rights form.
\copyrightyear{2022}
\setcopyright{acmcopyright}\acmConference[ICCAD '22]{IEEE/ACM International
Conference on Computer-Aided Design}{October 30-November 3, 2022}{San Diego, CA, USA}
\acmBooktitle{IEEE/ACM International Conference on Computer-Aided Design (ICCAD '22), October 30-November 3, 2022, San Diego, CA, USA}
\acmPrice{15.00}
\acmDOI{10.1145/3508352.3549370}
\acmISBN{978-1-4503-9217-4/22/10}


%%
%% Submission ID.
%% Use this when submitting an article to a sponsored event. You'll
%% receive a unique submission ID from the organizers
%% of the event, and this ID should be used as the parameter to this command.
% \acmSubmissionID{123-A56-BU3}

%%
%% The majority of ACM publications use numbered citations and
%% references.  The command \citestyle{authoryear} switches to the
%% "author year" style.
%%
%% If you are preparing content for an event
%% sponsored by ACM SIGGRAPH, you must use the "author year" style of
%% citations and references.
%% Uncommenting
%% the next command will enable that style.
%%\citestyle{acmauthoryear}

%%
%% end of the preamble, start of the body of the document source.
\begin{document}

%%
%% The "title" command has an optional parameter,
%% allowing the author to define a "short title" to be used in page headers.
% \title{The Name of the Title is Hope}

\title{HECTOR: A Multi-level Intermediate Representation for Hardware Synthesis Methodologies}


% \author{Ruifan Xu}
% \affiliation{
%   \institution{Peking University}
%   \city{Beijing}
%   \country{China}
%  }
% \email{xuruifan@pku.edu.cn}

% \author{Youwei Xiao}
% \affiliation{
%   \institution{Peking University}
%   \city{Beijing}
%   \country{China}
%  }
% \email{shallwe@pku.edu.cn}

% \author{Jin Luo}
% \affiliation{
%   \institution{Peking University}
%   \city{Beijing}
%   \country{China}
%  }
% \email{luo-jin@pku.edu.cn}

% \author{Yun Liang}
% \affiliation{
%   \institution{Beijing Advanced Innovation Center for Integrated Circuits}
%   \institution{Peking University}
%   \city{Beijing}
%   \country{China}
%  }
% \email{ericlyun@pku.edu.cn}

\author{Ruifan Xu$^1$, Youwei Xiao$^1$, Jin Luo$^1$, Yun Liang$^{1,2}$}
\affiliation{
    \institution{$^1$Peking University \\ 
    $^2$Beijing Advanced Innovation Center for Integrated Circuits}
}
\email{{xuruifan, shallwe, luo-jin, ericlyun}@pku.edu.cn}

%%
%% By default, the full list of authors will be used in the page
%% headers. Often, this list is too long, and will overlap
%% other information printed in the page headers. This command allows
%% the author to define a more concise list
%% of authors' names for this purpose.
% \renewcommand{\shortauthors}{Trovato and Tobin, et al.}

%%
%% The abstract is a short summary of the work to be presented in the
%% article.
% \begin{abstract}
%   A clear and well-documented \LaTeX\ document is presented as an
%   article formatted for publication by ACM in a conference proceedings
%   or journal publication. Based on the ``acmart'' document class, this
%   article presents and explains many of the common variations, as well
%   as many of the formatting elements an author may use in the
%   preparation of the documentation of their work.
% \end{abstract}

%%
%% The code below is generated by the tool at http://dl.acm.org/ccs.cfm.
%% Please copy and paste the code instead of the example below.
%%
\begin{CCSXML}
<ccs2012>
   <concept>
       <concept_id>10010583.10010682.10010689</concept_id>
       <concept_desc>Hardware~Hardware description languages and compilation</concept_desc>
       <concept_significance>500</concept_significance>
       </concept>
 </ccs2012>
\end{CCSXML}

\ccsdesc[500]{Hardware~Hardware description languages and compilation}


\keywords{Intermediate Representation, Hardware Synthesis}

%\ccsdesc[500]{Computer systems organization~Embedded systems}
%\ccsdesc[300]{Computer systems organization~Redundancy}
%\ccsdesc{Computer systems organization~Robotics}
%\ccsdesc[100]{Networks~Network reliability}

%%
%% Keywords. The author(s) should pick words that accurately describe
%% the work being presented. Separate the keywords with commas.
% \keywords{HLS}

%% A "teaser" image appears between the author and affiliation
%% information and the body of the document, and typically spans the
%% page.
% \begin{teaserfigure}
%   \includegraphics[width=\textwidth]{sampleteaser}
%   \caption{Seattle Mariners at Spring Training, 2010.}
%   \Description{Enjoying the baseball game from the third-base
%   seats. Ichiro Suzuki preparing to bat.}
%   \label{fig:teaser}
% \end{teaserfigure}

\definecolor{LightGray}{gray}{0.9}
\definecolor{myGreen}{rgb}{0.218, 0.9, 0.4}

\newcommand \dsir{\texttt{ToR}\xspace}
\newcommand \fsm{\texttt{HEC}\xspace}
\newcommand \tor{\texttt{ToR}\xspace}
\newcommand \hec{\texttt{HEC}\xspace}

\newcommand*{\rom}[1]{\uppercase\expandafter{\romannumeral #1\relax}}


\newlist{rules}{enumerate}{1}

\setlist[rules, 1]
{label=\Roman{rulesi}., %1., 2., 3., ...
leftmargin=5pt+\parindent,
rightmargin=5pt
}

\input{Abstract}
%\keywords{High-level Synthesis, Intermediate Representation, Dynamic Scheduling, Static Scheduling}
%%
%% This command processes the author and affiliation and title
%% information and builds the first part of the formatted document.
\maketitle

\section{Introduction}
Power efficiency has become a critical design factor with the continual demands of emerging applications including machine learning and scientific computation. As alternatives to general purpose CPUs and GPUs , customized hardware accelerators~\cite{extensor, fpdnn, bigdata_fpga} like ASICs and FPGAs provide good energy efficiency and performance. However, hardware design is still a big challenge. Hardware description languages (HDLs) including Verilog~\cite{verilog} and VHDL~\cite{vhdl} are frequently used in industry for hardware design. These HDLs adopt a low level of abstraction including wires, registers, and gates, also known as register transfer level (RTL). However, HDLs' productivity is seriously hampered by this low level of abstraction, which makes it difficult for hardware designers.

Both high-level synthesis (HLS) and hardware generators have been proposed to improve the productivity of hardware design. HLS is a general method for automatically generating hardware from a behavioral description written in high-level languages such as C, C++, or OpenCL~\cite{Dynamatic_Josipovi, uIR_Sharifian, flexcl}. HLS offers a promising way to design hardware in a high-level abstraction, which provides opportunities for software engineers without hardware experience to design hardware quickly. However, it can lead to bad performance and high resource utilization due to the lack of domain knowledge about the computation and architecture~\cite{hls_liang}. On the other hand, hardware generator makes use of domain knowledge to improve productivity. These generators often focus on a specific domain, such as streaming~\cite{aetherling,parallel_pattern} and spatial accelerators~\cite{tensorlib,ems}. With the domain knowledge of the target applications, these approaches use specific languages and synthesis flows to generate hardware with good performance. Although these approaches provide high level abstraction, they can only support a specific class of applications.

%Existing HLS tools~\cite{VivadoHls, legup, bambu, soda} don't take readability and flexibility into account, which makes the generating process opaque and cryptic.} 

% \hl{this paragraph gives an overview of the systhesis flow for each technique.} And HLS compilers also implement a series of transformations that makes up an end-to-end flow from high-level language to HDL program. 
Furthermore, different synthesis tools often adopt different synthesis methodologies. HLS compilers convert the high-level description to hardware implementation in three steps: allocation, scheduling, and binding. The scheduling step, which determines the cycles for each operation, can be implemented using different algorithms including static, dynamic, and hybrid scheduling~\cite{Dynamatic_Josipovi, cheng_dss}. Hardware generators perform sophisticated architectural transformations to improve performance and resource consumption. In addition to the domain specific optimizations, hardware generators are often guided with specific hardware templates to generate hardware implementation for different applications~\cite{tensorlib, aetherling}. %\hl{As for hardware generators, a domain-specific language (DSL) is often proposed to describe computation and achieve inner optimizations. Each DSL compiler must implement a complete flow that can generate HDL in order to convert DSL into hardware, resulting in numerous redundant implementations.}

% \hl{this paragraph should introduce the flows in details such as IRs.}
% \hl{Adopting fully separate synthesis flows makes it difficult to combine different methods. HLS offers a convenient way to design hardware, but it can only achieve conservative parallelism due to the omission of computation characteristics. In addition, hardware generators can provide efficient hardware with domain-specific optimizations, making it appropriate for stressful circumstances. The combination of different methodologies is a superior option for achieving good performance while maintaining enough productivity. However, programmers can only connect the generated hardware at the RTL level to combine different methods.}
% \hl{this is not a good reason.}

%The spatial accelerator including PE and interconnection is finally generated by different hardware templates.

Both synthesis flows require substantial engineering efforts due to RTL’s weak expressivity and low level of abstraction. HLS tools take high-level languages as inputs and use the compiler infrastructure. For example, Vitis HLS~\cite{VivadoHls} adopts LLVM IR~\cite{lattner2004llvm} as its internal representation, which is a widely used software compilation IR. After that, additional passes are needed to convert to hardware semantics. As for hardware generators, domain specific languages (DSL) and optimizations are often used for description and optimization. For example, Tensorlib~\cite{tensorlib} selects tensor algebra as its representation, and maps it into the hardware PE array using space time transformation.  However, these different methodologies share some similarities in the intermediate representations, such as control logic and hardware generation from a high-level abstraction. Therefore, it is possible to have a unified IR for different synthesis approaches so that users can easily design new hardware synthesis techniques based on the same infrastructure and explore different methodologies.


% \hl{here, we should talk about how difficult is for each approach, it is better use detailed IRs and tool chains as examples, and emphasize that they are different}. HLS tools take high-level languages as input and gain enough optimizations from the compiler infrastructure. Therefore, the intermediate representation (IR) of the language is also adopted by these HLS tools. As for hardware generators, a domain-specific language (DSL) is frequently proposed to describe computation and achieve optimizations. All these synthesis methodologies need a complete flow that can generate synthesizable hardware, resulting in numerous redundant implementations. \hl{However, these flows share some similarities in the intermediate representations, then give some insights. Then, it will be great to have unified IR for different synthesis approaches so that}.

%To avoid duplicated implementations and to explore better combinations of multiple approaches, a flexible intermediate representation is needed.

In this paper, we propose Hector, a two-level IR providing a unified description for different hardware synthesis methodologies with expressivity and flexibility. The high-level IR (TOpological Representation) binds computations with a control graph annotated with timing information, while the low-level IR (Hierarchical Elastic Componenent) provides a concise way to describe various hardware components and elastic interconnections among them using customizable primitives. Both IRs provide a uniform representation of the control logic with various manners, but at different abstraction levels. The IRs in Hector are converted to synthesizable RTL programs through a series of transformations including time graph transformation, lowering pass, and RTL generation. The two-level IR and all the transformations are built on the multi-level compiler infrastructure (MLIR)~\cite{mlir}. 

The main contributions of this work are as follows:
\begin{itemize}
\item We propose Hector, a unified IR and synthesis framework supporting multiple hardware synthesis methodologies including high-level synthesis and hardware generator. 
\item We propose a two-level IR that is general enough for different hardware synthesis methdologies. 
\item We implement a compiler that lowers the IRs to synthesizable RTL through a series of analysis  and optimizations.

%Hector provides the ability to design hardware with multiple synthesis approaches. 

%\item We establish several transformation passes to lower from \tor to \hec considering hardware with different behaviors.
% \item We use Hector to implement three hardware synthesis methods including high-level synthesis, systolic array generator and an accelerator for sparse linear algebra, which further demonstrates Hector's convenience and extensiblity.
\end{itemize}

To demonstrate the expressivity and versatility, we implement three synthesis approaches based on Hector: a high-level synthesis (HLS) tool, a systolic array generator, and an accelerator component for sparse linear algebra. The experimental results show that the hardware generated by Hector's hybrid HLS scheduling improves the performance by 29$\%$ on average compared to that generated by the state-of-the-art HLS tool. The other two approaches outperform HLS
implementations in performance and productivity. \textit{Hector is open source at GitHub (\url{https://github.com/pku-liang/Hector})}. 
% \hl{please give the detailed experiment data.}
\section{Background}

\subsection{High Level Synthesis}
HLS offers a promising way to design hardware at a high-level abstraction such as C-like language, which can release users from designing at register transfer level. HLS compilers~\cite{VivadoHls, legup, Dynamatic_Josipovi, uIR_Sharifian} then automatically convert the high-level language to HDLs by three major processes: allocation, scheduling, and binding. Allocation places all compute units on the datapath, and binding determines where each operation executes.

% \hl{first gives an overview of high level synthesis.}
\begin{figure}[t]
  \centering
  \includegraphics[width=\linewidth]{figure/background.pdf}
  \caption{Static and dynamic scheduling in high level synthesis.}
  \label{fig:background}
  \vspace{-9pt}
\end{figure}

Scheduling is the most important step in modern HLS tools. There are two different scheduling methodologies: static and dynamic, as shown in \autoref{fig:background}. The most common flow of static approach consists of 4 steps~\cite{intro-hls}: software compilation, static scheduling, allocating \& binding, and finite state machine (FSM) construction, as shown in \autoref{fig:background}(a). Compilation transforms the high-level program into a software intermediate representation. Then, static scheduling algorithm determines the execution time for operations considering dependencies and resource constraints. For example, the add and shift-right operations are placed in the first cycle (\verb|C0|). Allocating \& binding algorithms explore the opportunity for resource sharing, allocate necessary resources, and bind operations to them. Finally, a finite state machine is constructed for RTL generation. 

Static scheduling~\cite{rl_hls, cdl_hls, thread_weaving} is very effective for statically predictable programs such as perfect loops.  Pipelining is a key optimization technique to exploit parallelism among multiple loop iterations. In the sequential execution, one iteration can only begin after the previous iteration is finished. Pipelining allows iterations to overlap which improves throughput and gain resource sharing. The distance between two adjacent iterations is called initial interval (II)~\cite{}, which is the indicator of throughput.

% \hl{it seems that dynamic scheduling is not used in the following parts.}
Dynamic scheduling generates a dataflow circuit by leveraging elastic units~\cite{elastic_circuit}, such as merge, branch, etc., as shown in \autoref{fig:background}(b). All data signals in dataflow circuits are accompanied by handshake signals, which are \textit{valid} and \textit{ready}, in opposite directions, indicating the availability of the next data from the source unit and the readiness of the target unit to accept it, respectively. Static scheduling makes conservative assumptions for unresolvable dependencies at compile time. Dynamic scheduling overcomes this inefficiency by postponing the scheduling decisions until run time, However, dynamic scheduling suffers from high resource consumption. Hybrid scheduling combines static and dynamic schedulings~\cite{cheng_dss, find_islands}.

\subsection{Intermediate Representation}
Intermediate representation is an important abstraction to simplify the design of compilation and synthesis flows. Compiler infrastructures like LLVM~\cite{lattner2004llvm} and GCC~\cite{novillo2003tree} apply machine-independent optimizations and generate code for different target architectures based on IRs. For example, Static single assignment (SSA) form~\cite{StaticSingleAssignment} and control dataflow graph (CDFG) representation~\cite{cdfg_allen} are widely used in compiler optimizations and static analysis. Many HLS tools~\cite{VivadoHls, legup, Dynamatic_Josipovi,scalehls,scalehls} adopt LLVM IR as their internal representation. However, LLVM is a pure software IR that doesn't contain any hardware information. \cite{LLHD_Schuiki, firrtl, ahir,uIR_Sharifian, Calyx_Rachit} present hardware IRs that provide low-level abstraction to support hardware design. \hl{let us add liqiang's TENET paper into the list.} 

%Calyx~\cite{Calyx_Rachit} combines software-like control flow and structural representation to describe architectural design.
%Hardware IR are used in hardware synthesis flows.

%Halide~\cite{Halide} aims at explore parallelism and locality of software domains like image processing.
% Relay IR~\cite{relay} proposes a differentiable language for expressing machine learning models. 

MLIR~\cite{mlir} is a novel compiler infrastructure that provides powerful scalability and modularity. MLIR greatly facilitates the implementation of various IRs and transformations among them. All IRs in MLIR follow the SSA form and an explicit type system. \emph{Dialect} in MLIR is a hierarchical structure template allowing basic optimizations in MLIR to be reused. This hierarchy supplies an expressive representation, which makes it easy to implement a flexible IR. 

% \section{Background}
% \subsection{HLS Scheduling}
% Traditional HLS tools like Vivado HLS\cite{VivadoHls} and LegUp\cite{legup} adopt static scheduling like SDC algorithm\cite{cong_sdc} that decides certain clock cycle in the compile time, allocates each operation to a state and uses Finite State Machine (FSM) to deal with control flow. Therefore, hardware implementation works in a sequential way. To increase parallelism, explicit pipeline optimizations like module scheduling~\cite{rau_module_scheduling} are adopted. When pipelining the program, static scheduling needs to make a conservative scheduling at the compile time that reserves enough time to ensure all dependence and memory conflicts cannot be violated. In a word, if the program has complicated control flow and the branches are imbalanced, these HLS tools may choose a big Initial Interval (II) which leads to a bad performance. One benefit of static scheduling is that we can analysis whether every pair of operations can share the same resource, and gain enough resource sharing.

% To overcome the conservative pipeline in the compile time, Dynamatic~\cite{Dynamatic_Josipovi} adopts elastic circuits techniques\cite{elastic_circuit} in dynamic scheduling. Each operation in CDFG is wrapped to a pipeline-style compute unit, and the control flow is converted to the combination of many elastic components such as fork, merge, branch and select. As for memory conflicts, their work adopts a Load Store Queue (LSQ)~\cite{load_store_queue} to dispute load and store operations. Every pair of operations that has dependence in CDFG is communicating through a handshake protocol\cite{handshake}. By the use of these signals, we can satisfy the dependence during the run-time and overcome the shortcoming of conservative selection. Lacking enough information at the compile time, resource sharing between different operations cannot be obtained. To get a trade-off between performance and resource utilization, DSS~\cite{cheng_dss} first combines these two scheduling method in functional level by manually determining whether each function uses dynamic scheduling. The static parts are then wrapped with handshake protocol, which are directly connected to dynamic parts at RTL level.

% \subsection{Intermediate Representation in HLS}
% Intermediate representation is an important abstraction that represents a model of software or hardware, and can be applied on compiler transformation for analysis and optimizations. As a powerful software compiler infrastructure, LLVM is often used in HLS tools like Vivado HLS and LegUp. Because LLVM IR doesn't aim at hardware, some inner representations are extra implemented in those tools. Existing HLS tools don't consider about readability and flexibility, for which these inner representations are not exposed to users.

% \subsection{MLIR Infrastructure}
% MLIR\cite{MLIR} is a novel compiler infrastructure which provides powerful scalability and modularity, and is particularly suitable for multi-level IR transformation. The \emph{Dialect} in MLIR is a set of operations, attributes and types, which can be regarded as a flexible IR. All the IRs in MLIR obey Static Single Assignment (SSA) form~\cite{StaticSingleAssignment} and explicit type system. MLIR provides several built-in Dialects, such as \emph{Affine}, \emph{SCF} and \emph{Standard}. \emph{Affine} provides a abstraction for affine operations like affine.for, which can be used to polyhedral compilation techniques. \emph{SCF} describes static control flow in a higher-level abstraction keeping branch and loop structure, which is helpful for analyzing and optimizing. \emph{Standard} is a collection of basic operation such as addition, comparison and function call. These dialects can be used to describe a program at the same time, which is also the feature provided by MLIR, called progressive lowering. The program described in \emph{Affine}, \emph{SCF} and \emph{Std} has similar functionality and representation with high-level language.

% \subsection{Motivation}
% Most of HLS tools only adopts static scheduling, while several tools such as Dynamatic and SOFF\cite{SOFF} only considers about dynamic scheduling. DSS proposes a mixed scheduling for performance and resource trade-offs. However, the mixed scheduling method that DSS designs is just the simple combination of dynamic and static parts. The method proposed by DSS can be seen as a function-level mixed scheduling approach, which needs an elaborate design to partition an application into multiple kernels. Then, dynamic and static parts are separately sent to Dynamatic and Vivado HLS, which are finally lowered to VHDL. This scheduling method isolates two kinds of scheduling algorithms, which makes the hardware generation separate and independent.

% High-level synthesis substantially reduces the difficulty on hardware design by automatically converting high-level programs to HDLs. But traditional HLS frameworks lack explicit intermediate representation, which makes it hard to inspect and interfere during the whole HLS process. It is the only way to manipulate the hardware generation that users manually insert proper directives like loop unroll and pipeline. Without the information of each step, we can only get the final HDL implementation which is not friendly for understanding and further optimization. In another word, the area and performance of final hardware can only count on HLS compilers, which may generate unpredictable hardware with different behaviors from what we expect. Different HLS tools support different kinds of directives, and may use different scheduling algorithm which cannot be interfered or selected. Explicit IR can provide significant convenience to generate predictable hardware especially during the scheduling and binding phases. For example, users can select different scheduling algorithm provided by our flow, and they can even implement their own scheduling algorithm to generate better hardware with more knowledge of their own applications and reuse the other part in our flow. What's more, some optimizations like resource sharing and memory partition are easily applied, which can also be embodied in this IR.

% So we are willing to implement a HLS framework based on multi-level IR paradigm that is suitable for representing hardware from different abstraction levels and that is also easy for compiler to optimize. To make use of both static and dynamic scheduling, the IR we design should also be capable of representing both dynamic and static scheduling. An IR designed for HLS lessens the gap between high level programs and HDLs, and also provide the ability to debug and manipulate at different level.

% Hardware generated by Vivado HLS can only be interfered by manually inserting directives, while DSS only supports the schedule selection of each function. In contrast, the hardware can be interfered at different level like scheduling and binding strategies with the help of explicit IR.
% \section{HECTOR Framework}
%\subsection{MLIR Infrastructure}
% \hl{let us briefly discuss this in flow.}

% Hector is built on MLIR\cite{MLIR}, a novel compiler infrastructure providing strong scalability and modularity, and is particularly suitable for multi-level IR transformation. The \emph{Dialect} in MLIR is a set of operations, attributes and types, which can be regarded as a flexible IR. All the IRs in MLIR obey Static Single Assignment (SSA) form \cite{StaticSingleAssignment} and explicit type system. MLIR provides several built-in Dialects, such as \texttt{SCF} and \texttt{Standard}. \texttt{SCF} describes static control flow in a higher-level abstraction keeping branch and loop structure, which is helpful for analyzing and optimizing. \texttt{Standard} is a collection of basic operation such as addition, comparison and function call. These dialects can be used to describe a program at the same time, which is also the feature provided by MLIR, called progressive lowering. The program described in \texttt{SCF} and \texttt{Standard} has similar functionality and representation with high-level language.

%\subsection{Overall flow}

% \begin{figure}[t]
%   \centering
%   \includegraphics[width=\linewidth]{figure/overview.pdf}
%   \caption{Overview of HLS flow with Hector framework}
%   \label{fig:flow} \vspace{-0.3cm}
% \end{figure}

% \autoref{fig:flow} presents the overall flow of Hector. The framework is built upon MLIR  \cite{MLIR} ecosystem, a highly customizable compiler infrastructure, which greatly facilitates the implementation of various IRs and transformations among them. The input to Hector is standard C or C++ program. We adopt open-source tools Clang \cite{lattner2008llvm} and Polygeist \cite{polygeist}, which lower C/C++ description to MLIR system in a two-step manner. Hector then performs an efficient algorithm (as detailed in \autoref{sec:algorithm}) to partition the design into submodules, synthesizing them progressively in a hybrid manner, and generate the Chisel description at the register transfer level (RTL). 

%Our framework not only supports versatile scheduling techniques, but also provides lots of customization and extension opportunities thanks to the implementation of algorithms upon MLIR infrastructure. 

%Hector applies compilation optimizations, such as constant propagation, on \texttt{scf} dialect, an existing software IR with high-level control flow semantics in MLIR system. The optimized \texttt{scf} code is transformed into a raw \tor description through operation rewriting and topology initiation. 

% Hector is essentially a two-level IR system equipped with a set of optimization and transformation passes. We design \tor IR at a high level and \hec IR at a low level, as detailed in \autoref{sec:ir}. \tor (Topology Rail) IR features a directed graph, where high-level operations are bound on nodes or edges. Compared with pure software IRs, such as LLVM IR \cite{lattner2008llvm}, \tor is much more capable for hardware design because the graph carries rich hardware semantics. \hec (Hybrid Elastic Component) IR is designed at the near-RTL level. It describes each component in either finite-state machine (FSM) or pipelining manner, and integrates them using elastic units. The hybrid nature of \hec greatly facilitates the description of hardware generated by versatile HLS methodologies, while existing hardware IRs either support only a specific type of hardware (\cite{Calyx_Rachit} only uses FSM to describe hardware control, and \cite{uIR_Sharifian} only describes dataflow circuits, for example) or \cite{LLHD_Schuiki} describes digital circuits at a much lower level not suitable for HLS flow.

%and bridges the semantic gap between Von Neumann programming model and fine-grained parallel hardware. 


% After execution, Hector would deliver the produced RTL description with appropriate IP configurations to the commercial software Vivado for logic synthesis and physical design, which yields bitstream results together with analysis reports.

% \autoref{fig:flow} depicts how a high-level program is transformed to hardware in Hector. Hector first compiles C/C++ programs to clang AST, which Polygeist\cite{polygeist} traverses to create \texttt{SCF} + \texttt{Standard}. Some essential optimizations like constant propagation are applied, and a raw \dsir is built by rewriting operations and initializing the schedule topology. Hector can construct a hybrid scheduling circuit and then adopt a static scheduling techniques on static modules, both of which can be simply implemented on the topology structure of \dsir. The schedule results are then lowered to \fsm composed of different components using typical static and dynamic HLS synthesis processes. \fsm components are directly transformed to Chisel\cite{chisel} programs that describes the hardware using some pre-defined Chisel templates. The simulation provided by Chisel can be used to validate the functionality and performance, and Chisel can generate results written in Verilog which fits in logic synthesis and physical design.

% We design two a two-level IR, \dsir and \fsm, which are hardware descriptions at different abstraction level. \dsir{} provides a hardware description IR with a well-designed type system and operations, which represents static and dynamic scheduled compound design in a hierarchical description. \fsm{} reduces the gap between two scheduling strategies, and is also a hardware-like structure which can be easily translated to hardware implementation like HDL. These two dialects are introduced in \autoref{sec:dialect}. 

% Our proposed flow is built on these two \emph{Dialects} and MLIR build-in high-level dialects. The input dialects includes \emph{SCF} and \emph{Std}. These dialects together describe a high-level program, which reserves the structures of branch and loop. Compared with low-level representation CDFG, the structures of branch and loop are friendly to understand and scheduling. Furthermore, \emph{Affine} defines a collection of affine operation, which can make use of polyhedral model analysis to achieve loop transformation and memory partition. Optimization directives like loop unroll and memory partition can also be inserted into this dialect. Then, static scheduling is applied to the high-level program, where SDC algorithm is provided in this flow. This framework can also support other scheduling algorithm that is implemented by users. Static scheduling process will generate fully static modules represented in \dsir{}. Then, static modules are transformed to mixed modules by the a partition algorithm, where mixed modules are the combination of static and dynamic modules. The mixed modules are then split into static and dynamic modules, both of which have further transformation passes. After that, both of modules are lowered to \fsm{}. Having done some low level transformations in \fsm{}, the programs are then converted to Chisel\cite{chisel}. Detailed design and the reason why partition algorithm is applied to the static modules are introduced in the following sections.

\begin{figure*}[htbp!]
  \centering
  \includegraphics[width=\textwidth]{figure/tor.pdf}
  \caption{The design of \tor IR. \tor IR consists of topology and functional operations. In (a), topology describes the time graph with supplementary information on edges, where all edges and nodes are set "static" as default. In (e), operations are bound on the graph according to the syntax in (d). The binding situation determines the types of nodes on the time graph, as shown in (b). (c) shows the restrictions the node types set on the graph structure.}
  \label{fig:tor}
  \vspace{-9pt}
\end{figure*}

\section{Intermediate Representation}\label{sec:ir}
% \hl{let us first write a paragraph about the novelty of intermediate representation, then talks about its implementation and what are the proposed IRs used for? What are the relations between different IRs.}
In this section, we present the details of the IRs , \tor (TOpological Representation) and \hec (Hierarchical Elastic Component). 

\subsection{Overview}
% \hl{it is not clear how these IRs can be used for different hardware synthesis approaches.}

Hector is a two-level IR system, where \tor is at the high level and \hec IR is at the low level. \tor IR combines a software-like control flow and the schedule information of each operation. \hec IR proposes an allocate-assign mechanism to explicitly describe the relationship between computation and compute units. Both IRs provide a uniform representation of the control logic with various manners. The main difference between the two IRs is that \tor describes when the operation begins, while \hec describes where it takes place. Hector is built on a novel compiler infrastructure, MLIR~\cite{mlir}, that provides powerful scalability and modularity.

Because of the two-level representation, it is possible to use multiple hardware synthesis methodologies. \tor is capable of providing a high-level abstraction of the schedule information. Both dynamic and static scheduling methods in HLS can be easily captured by \tor IR. The allocation of hardware resources and the binding of operations can then be described in \hec IR. The entire HLS procedure can be obtained from several lowering transformations based on Hector IR. Hector also can describe hardware at the architectural level, which is a necessary task for hardware generators. The allocate-assign method in \hec makes it easier to describe interconnection at the low level, while high-level scheduling transformations can be applied to \tor IR. %\hl{this paragraph is not sufficient, we need more details here.}



\subsection{ToR IR} \label{sec:tor}
The software IR such as LLVM IR lacks of hardware semantics. The idea of the high-level IR is to make it closer to hardware by providing a directed graph that carries control flow and timing information and binding software operations to elements of the graph. \tor is composed of two parts, topology and functional operations, both of which are introduced below. 

Topology describes a time graph, which is a directed graph describing control flow and timing information. Topology includes a \verb|tor.topo (x to y)| operation, which indicates the source node and sink node numbered \verb|x| and \verb|y|, respectively. The \verb|tor.from| operations inside \verb|tor.topo| specify edges of the time graph. Attributes such as \verb|"seq:x"| add supplementary information to the time graph. There are four types of nodes in the time graph. Normal node leads a sequential edge, where \verb|"seq:2"| indicates that the bound operation occupies two cycles. Call node leads a call edge, where \verb|"call"| indicates a state that stalls until the callee finishes. If node leads two branches, and loop node leads a loop body and loop back edge. The type of each node on the time graph is determined by the operation binding. The combination of these four node types is capable of describing the schedule at a high-level.

Three scheduling manners: \textbf{static}, \textbf{pipeline}, and \textbf{dynamic}, are supported in \tor. Pipelining is a key optimization technique to improve throughput, which serves as the determinant of the overall performance. \tor IR supports pipelining by aligning branches of all \verb|if| operations and adding \verb|pipeline| and \verb|II| attributes to submodules. Topology also supports dynamic behavior that resolves conflicts in run-time. Stalling occurs only when the conflict occurs, avoiding the conservative assumption of static behaviors. This unified representation makes it easier to transform among different behaviors. 

Functional operations present the algorithmic specification with high-level control flow semantics (e.g., if, for, and while). It binds each operation to some element of the time graph, either a node or an edge. To be specific, general operations (computation, memory access, function call) are bound on edges, while if/loop operations are bound on nodes. Functional operations specify functionality and binding according to the syntax in \autoref{fig:tor}(d). 

\autoref{fig:tor} illustrates an example of \tor IR. The time graph in (a) contains a loop, which is composed of two branches $1 \to 2 \to 3$ and $1 \to 4$. There is also a function call on the edge $2\to 3$, and the loop exits at the edge $0\to 7$. As shown in (b)(c), four types of node vary in the graph structure. (e) shows the functional operations which is bound on the time graph. For example, the \verb|tor.for| operation is bound on node \verb|0|, and the \verb|tor.muli| is bound on edge \verb|(1 to 2)|. As shown in \autoref{fig:compile}(a), there can be multiple manners in the same time graph,  which are attached to the time graph as attributes like \verb|"seq:2"| and \verb|"dynamic"|. The time graph with multiple manners is called a hybrid time graph. 

\subsection{HEC IR} \label{sec:hec}
Here, we propose \hec IR, which describes hardware with different manners in a unified allocate-assign mechanism, as depicted in \autoref{fig:hec}(a)(b). Allocation explicitly defines all function units and sub-modules on the datapath, and the signals of these units are determined through assignments in a specific situation. That allocate-assign mechanism omits the insertion of multiplexer, which simplifies the assignment of signals.

Compared with \tor, \hec works at a level much closer to hardware. It explicitly describes the resource usage (including registers, on-chip memory, and computation units). Corresponding to the different behaviors in \tor: \textbf{static}, \textbf{pipeline} and \textbf{dynamic}, a \hec design is composed of three types of components matching their different manners as follows. 

\begin{figure}[htbp!]
  \centering
  \includegraphics[width=\linewidth]{figure/hec.pdf}
  \caption{Allocate-assign mechanism and three styles of components in \hec. 
   (a) shows an example for unit allocation, where hec.primitive allocates a "muli" unit named "m" with three ports, and hec.wire declares a wire, which is commonly used in HDLs. (b) shows the usage of combinational operations and hec.assign for signal delivery among ports and wires. (c)-(e) presents three component styles. }
  \label{fig:hec}
  \vspace{-9pt}
\end{figure}

\textbf{STG-style component}. 
\hec describes a static module in a state transition graph (STG) style. The \verb|hec.stateset| operation defines a set of states, as shown in \autoref{fig:hec}(c). Inside each state \verb|@sx|, \verb|hec.assign| operations specify the signal delivery situation among allocated resources. Such representation naturally supports fine-grained parallelism. There is also a \verb|tor.transition| operation in every state, specifying which state the control is transferred to, either unconditionally or based on guard signals. Based on the allocate-assign mechanism, it is convenient to describe resource sharing in an STG-style component by simply feeding signals into the shared resources (either registers or compute units) inside different states.

\textbf{Pipeline-style component}.
The component for a pipelined module is described in a multi-stage style. \hec explicitly presents all pipeline stages by \verb|hec.stage| operations inside the \verb|hec.stageset| operation, as shown in \autoref{fig:hec}(d). For each value, \hec allocates a register to carry it for each stage between its definition and the latest use. This is because there may be time overlap between the lifetimes of the same value in consecutive executions. Signal delivery description is similar to that of STG-style components, while resource sharing additionally needs to consider the initial interval (II) for conflict avoidance~\cite{sdc-modulo-13, dai_modulo_scheduling}. We also define operation \verb|hec.deliver| to specify inter-iteration data delivery.

\textbf{Handshake-style component}.
\hec describes the dynamic submodule in handshake style. In order to simplify the description, extra signals in the handshake protocol, such as valid and ready, are hidden in this kind of component. With elastic units, such as \verb|branch| and \verb|fork|, predefined as primitives in \hec, it is easy to describe the interconnections among components with the allocate-assign mechanism inside the \verb|hec.graph| operation, as shown in \autoref{fig:hec}(e). 

\subsection{IR Implementation} 
Both IRs are implemented as \emph{dialects} in MLIR infrastructure. MLIR provides a generic form of operations, which can be customized for new languages. Besides basic operands and results, operations may have \emph{attributes} and \emph{regions}. \autoref{fig:tor}(d) shows the functional syntax of operations in \tor. There are two \emph{regions} in \verb|tor.if| representing two branches separately. The nested region is naturally supported in MLIR, which enables the definition of \verb|stateset| and \verb|stageset| in \autoref{fig:hec}(c)(d). The different behaviors of schedule, the binding of operations and the instantiation of sub-module are all implemented as \emph{attributes} including literal and symbolic reference.

The two-level representation makes it easy to implement different synthesis methods. \tor IR provide a high-level abstraction of schedule information, and different scheduling approach can be easily obtained by transforming on the time graph. The explicit representation of allocation in \hec IR brings the opportunity of resource sharing, which significantly reduces the resource consumption. The allocation of sub-modules forms a hierarchical representation of hardware, which is capable of describing architectural design. The allocate-assign mechanism also simplifies the definition of interconnection between different modules. Therefore, both HLS and hardware generators can be easily implemented in Hector IR.

% 1. how to build hls and hardware generators using two IRs.
% 2. how to use MLIR to build the two IRs is not clear.
\begin{figure*}[htbp!]
  \centering
  \includegraphics[width=\textwidth]{figure/compile.pdf}
  \caption{An overall example for compiling Hector to hardware. The original module with hybrid manners is split into two modules with dynamic and pipeline manners separately, and the connection between them is represented as a function call. Through resource and register sharing, pure \tor modules are transformed to \hec components containing compute units and registers. RTL implementations are generated by mapping ports and constructing controllers with the same functionality.}
  \label{fig:compile}
  \vspace{-9pt}
\end{figure*}

\section{Compiling IRs to Hardware}
\label{sec:compile}
%With the help of MLIR infrastructure, basic optimizations like dead code elimination (DCE) and common sub-expression elimination (CSE) can be easily implemented. Apart from these basic optimizations, hardware generation in Hector should go through these transformation passes:

To compile IR to hardware, Hector involves the following passes, Hybrid Time Graph Transformation (\autoref{sec:split}), Lowering from \tor to \hec (\autoref{sec:lower}), and RTL Generation (\autoref{sec:rtl}). \autoref{fig:compile} illustrates the main steps from \hec to RTL implementation written in Chisel. 
% \hl{can we change figure 4 to explain all the three steps ?}

The first pass converts the hybrid time graph in \tor into multiple pure time graphs by outlining connected sub-graphs. The second pass produces \hec programs with the same functionality of \tor through pipeline partition, control logic generation, and resource sharing. The third pass generates Chisel programs by implementing hardware controller and inserting necessary signals.

% \hl{we need a figure to explain this pass: inputs, outputs, transformations, etc. }

% \subsection{Verify Functionality}
% \label{sec:verify}
% The high-level topology enables early verification \hl{verification of what ? this needs examples} at a high level of abstraction. Some basic requirements of hardware design can be ensured by traversing the time graph. 

% \begin{itemize}
%     \item \textbf{Latency}: Because one variable cannot be used until the computation is complete, each path between its definition and usage should be shorter than the operation latency.
%     \item \textbf{Timing}: The total length of each combinational path cannot exceed the clock period, which can be examined by scanning each time edge.
%     \item \textbf{Memory Port}: The number of accesses in a single edge cannot be greater than the product of port limitation and edge length due to memory port limitations.
% \end{itemize}

\subsection{Hybrid Time Graph Transformation}
\label{sec:split}
\tor's semantics allows for the existence of multiple manners in the same time edge. However, because different types of hardware cannot communicate directly with one another, this hybrid time graph is unsuitable for subsequent transformations. This pass converts the hybrid time graph into multiple pure time graphs where all edges and nodes have the same manner. The connected sub-graph and related computation are defined as an external module, and communication between the outer and inner modules is accomplished through a function call. 

We calculate all live-in variables and live-out variables of connected sub-graph. Live-in variables are the arguments of new function, and that function return the values of all live-out variables. As shown in \autoref{fig:compile}(b), the connected sub-graph of hybrid time graph is marked by red rectangle. That hybrid time graph is split into two sub-graphs: a dynamic module and a pipeline module. The inner pipeline module take \verb|%c| and \verb|%d| as inputs, and return the value of \verb|%f|. This pass ensures that only one type of hardware is considered in the following transformations. 
% \hl{this subsection needs more details.}


\subsection{Lowering \tor to \hec}
\label{sec:lower}
The time graph in static modules is converted to a low-level FSM with the same functionality. For pipeline \tor modules, this pass traverses the time graph and calculates the "distance" of each node from the source node. The "distance" is indicated by “\#cycle” attributes of every time edge composing the path from the source node to the current node. The time nodes with the same "distance" will be placed in one stage. Each stage corresponds to a clock cycle, on which the allocated time node starts the execution of the operations whose start time is set as the current node. 

In contrast to \tor, \hec specifies where the computation takes place and where the intermediate value is stored. This opens up the possibility of resource and register sharing, which allows disjoint computations and values to share the same unit. \tor fits well in such optimizations because of the high-level control flow and timing information of computation contained in the time graph.

\textbf{Resource and register sharing}. The resource sharing can only happen between two operations that never occur at the same time, so two disjoint operations mustn't overlap in the time graph. We build a conflict graph that shows all potential conflicts and uses a coloring strategy to assign compute units to each node. As for register sharing, we use a live range analysis on the time graph to calculate the live range of each variables. The value of a variable in its live range must be stored in a register unless it is used in time. The main difference between register sharing and resource sharing is that each edge in the conflict graph represents an overlapping live range. 
% \hl{in 4.1, we mention a liveness analysis, is it the same ?}

As for hardware with dynamic behavior, each pair of operations may happen at the same time, restricting resource sharing. As a result, the lower pass just considers the generation of control logic. The handshake protocol, which describes data transfer with valid and ready signals, is used to implement the dynamic approach. This protocol is hidden in the \hec IR, which further simplifies the description. 
% \hl{we do not introduce the elastic components before.}
Elastic components~\cite{elastic_circuit} are adopted to solve the synchronization of tokens in the hardware with dynamic behavior. For example, \verb|branch| transfers the data to one of the outputs based on a condition, and \verb|merge| accepts one of the inputs which is similar to a $\phi$-function in software IR. We calculate the live-in and live-out variables of each regions, and a template-based approach is adopted to generate data transfer of variables among different regions. After that, \verb|fork| is inserted to synchronize multiple uses of the same value, which is shown in \autoref{fig:compile}(c). Both \verb|addf| and \verb|subf| need the value of \verb|%a|, so a \verb|fork| component is inserted for that variable. The two outputs of \verb|fork| are fed into \verb|addf| and \verb|subf| separately.

\subsection{RTL Generation}
\label{sec:rtl}
\hec contains sufficient information to generate synthesizable RTL. Each component in \hec has explicit I/O port definitions, assignments of ports, a controller, and allocations of register and resource (including memory, compute units, and other sub-component). The RTL generation pass creates a Chisel \cite{chisel} program by mapping each component to a Chisel module with the same ports and assigning the corresponding value to these ports. Chisel will automatically insert the multiplexer, allowing for multiple assignments under different conditions. All of the built-in resources, such as memory and compute units are implemented as a pre-defined Chisel Module. 

\textbf{Controller Generation}. The STG-style component is controlled by a state transition graph that contains a state set and the transition between them, whereas the pipeline-style component is controlled by a collection of stages and matching II. By inserting some extra registers and wires, this pass builds a Chisel-style FSM for these two types of components with the same functionality in \tor. Due to the need to resolve conflicts at run-time, handshake-style components do not have a centralized controller. For each port, the handshake protocol is implemented as a Chisel class \texttt{DecoupledIO} with ready, valid, and data signals. The Chisel program for a pipeline component and a dynamic component is shown in \autoref{fig:compile}. The controller of the pipeline component is eliminated because the II equals one.

\textbf{Wrapper construction}. Due to the distinction between handshake-style component and the other two types of components, this pass additionally generates a wrapper that attaches the handshake protocol to naked components. Whether an STG-style (or pipeline-style) component is valid and ready is determined by the component's completion (or the validity of the initial stage and final stage).
% \section{Algorithm} \label{sec:algorithm}
Here, we develop a heuristic partitioning algorithm, which divides the original specification into submodules and determines each submodule to be scheduled, whether statically or dynamically. 
The insight of our heuristic is to apply dynamic scheduling only when necessary and partition the origin specification into as large submodules as possible, each of which will be statically scheduled and then integrated into the whole dataflow circuit. 
The algorithm is performed at the \dsir{} level, where the structure of the \tor graph facilitates the extraction of static submodules. 

% \begin{algorithm}
% \caption{2-step Partitioning}
% \label{alg:partition}
% \setstretch{0}
% \KwIn{Raw \tor description $r$\tor}
% \KwOut{Hybrid \tor description $h$\tor}
% set $\chi(op)=\mbox{S}$ for every operation $op$ \\
% \ForEach{inner-most loop $l$}{
%   \ForEach{$if \in l$ which guards any computation of $l$'s inter-iteration dependencies}{
%     set $\chi(if) = \mbox{D}$ \Comment{Rule \rom{1}}
%   }
%   \ForEach{memory op $y \in l$ with non-affine indices}{
%     set $\chi(y) = \mbox{D}$ \Comment{Rule \rom{2}}
%   }
% }
% \ForEach{loop/if $x$ which contains $y:\chi(y)=\mbox{D}$}{
%     set $\chi(x) = \mbox{D}$ \Comment{Rule \rom{3}} \\
% }

% extract the static path set $P$, and set submodules\\
% generate $h$\tor.
% \end{algorithm}

% Algorithm 1 consists of  two main steps: marking and partitioning, detailed as follows.
\subsection{Marking}

To determine the scheduling manner of operations, the marking step marks nodes and edges in \tor graph either static or dynamic. 
We observe that the performance of the whole design is mainly determined by the inner-most loops, which contain the the most executed operations. 
For these loops, dynamic scheduling has obvious performance advantages against the static one in the presence of undetermined inter-iteration dependencies, either irregular memory accesses or unpredictable control flows.

We propose the following rules to select the operations need to be dynamically scheduled.

\begin{rules}
    \item \textbf{If operations} \textit{guarding any computation of an inner-most loop's inter-iteration dependencies.} Because dynamic scheduling may alleviate the negative effects of such guarded inter-iteration dependencies.
    \item \textbf{Memory operations} \textit{with unresolvable indices inside an inner-most loop.} Because static scheduling has to assume worst-case dependencies among them, which limits performance.
    \item \textbf{Loop/if operations} \textit{containing any dynamic operation.} Because it is impractical to embed any elastic component in a statically scheduled circuit.
    % \item \textit{Set of memory operations which access the same memory but lie on distinct static paths.} Because it is not safe to schedule these memory operations locally without a global load-store queue (LSQ). Furthermore, the LSQ would stall the execution of static submodules, which effects their performance.
\end{rules}

% As depicted in \autoref{alg:partition}, Rule \rom{1}, \rom{2}, and \rom{3} are carried out at lines 3-4, 5-6, and 7-8, respectively.


\begin{figure}[htb!]
    \centering
    \includegraphics[width=\linewidth]{figure/algorithm.pdf}
    \caption{An example for the partitioning algorithm. In step 1, the elements of \tor graph get marked following Rules \rom{1}-\rom{3}, as depicted by code on the right. In step 2, the algorithm extracts static (blue) submodules from the \tor graph and generates corresponding tor.call operations in the dynamic submodule. The \tor prefix is omitted for simplicity.}
    \label{fig:algorithm}
\end{figure}


\subsection{Partitioning}
Our algorithm further partitions the design into submodules based on the marked \tor graph, which is shown as Step 2 in \autoref{fig:algorithm}. 
The regular structure of the topology facilitates the step, which only needs to extract the static submodules composed of consecutive static nodes and edges on the \tor graph. The algorithm allocates submodules with corresponding operations placed inside, and generates the hybrid \tor. The produced hybrid \tor contains one dynamic submodule and several static submodules. The dynamic submodule includes all dynamic operations as well as generated function calls to static submodules.
% After producing a hybrid scheduling solution, our proposed holistic framework synthesizes the design into an efficient circuit with good performance as well. We will demonstrate the algorithm's effectiveness in \autoref{sec:experiment}. \hl{we need to more details.}




% \section{Transformation Pass}

\tor and \hec IRs as well as transformation passes are built on top of the MLIR framework. To generate the final design, we implement the following passes in Hector: 
\begin{enumerate*}
\item Static scheduling pass: this pass performs static scheduling on input IR and generates raw \tor; \item Partition and re-scheduling pass: this pass transforms raw \tor to hybrid \tor by partitioning raw \tor into static and dynamic submodules. Users of Hector can implement their own partition algorithm with the scheduling information provided by raw \tor. Then each static submodule in hybrid \tor is re-scheduled independently; \item \tor lowering pass: this pass lowers the scheduled \tor to \hec; \item Chisel generation pass: this pass generates Chisel code from \hec description.
\end{enumerate*}


\textbf{Static scheduling pass}.
%SDC-based scheduling\cite{sdc-06} is a static scheduling algorithm that can handle various kinds of constraints and perform optimization against diverse objectives. The same method is adopted in modulo scheduling for pipeline synthesis\cite{sdc-modulo-13}. 
In our implementation, we use the SDC-based scheduling~\cite{sdc-modulo-13} as the default scheduling strategy to convert the front-end generated IR to static scheduled \dsir{}.

\textbf{Partitioning and re-scheduling pass}.
Section \ref{sec:algorithm} describes our algorithm to partition a raw \tor into static submodules. More sophisticated partition algorithms can be used in this pass as substitutions to our algorithm. After partitioning, each static submodule is re-scheduled using the same algorithm as static scheduling pass.

\textbf{\tor lowering pass}. 
Each submodule's \tor graph structure incorporates scheduling results and provides useful information such as the operation's live range and mutually exclusive operations. Such information is useful for subsequent transformations and optimizations like resource binding. It greatly eases the difficulty of designing \tor lowering pass. The \tor lowering pass allocates computation and storage units for each static submodule with the help of the \tor graph. For an unpipelined circuit, its control logic can be deduced directly from the time graph by simply making each node in the \tor graph as a control state and generating state transition logic based on the edges in the \tor graph. For pipelined circuits, the cycle number associated with each edge helps us divide the circuits into stages. As for dynamic submodules, its \tor graph inherently implies the corresponding elastic circuit. Elastic components can be inserted based on their explicit control statements. Each submodule is instantiated using \verb|hec.instance| and each allocated unit is instantiated using \verb|hec.primitive|. Their inputs and outputs are connected using \verb|hec.assign|, which forms a graph representation. A buffer insertion algorithm is then applied to avoid pipelining stall by traversing the graph in \hec. 

%We leave a sophisticated buffer insertion algorithm as the future work.

% The graph-based representation of handshake. Buffer insertion, we leave a sophisticated buffer insertion algorithm as the future work.

% \hl{dynamic optimizations like buffer insertion}

\textbf{Chisel generation pass}. 
% \hl{this paragraph is not clear.}
The allocate-assign mechanism of \hec is naturally close to HDLs at the register-transfer level (RTL). According to \hec IR, this pass generates Chisel descriptions to allocate resources and specifies the value delivery among ports and wires. 
 The behavior of STG components is controlled by a finite state machine (fsm). A shift register along with a fsm are instantiated to construct a pipeline structure. Each port in handshake component is encapsulated in class \texttt{DecoupledIO}, which hides the valid and ready signals inside.


% \fsm{} is closer to hardware, which explicitly defines all needed connections, registers and compute units. So it doesn't need many transformations that convert the components in \fsm{} to chisel program. All the primitives in \fsm{} are implemented as a chisel class, which includes arithmetic operations and elastic components. Chisel provides the support for handshake protocol, which is encapsulated in class \textbf{DecoupledIO}. Each \textbf{fsm.primitive} and \textbf{fsm.instance} is converted to a instance. Because the port name is stored in the result of operation, these port can be easily to access like \textbf{adder.result}. \textbf{stateset} region is translated to the finite state machine in Chisel which has the same functionality except the go insertion in the initial state. The multiplexer will be automatically inserted by Chisel, which makes it possible to multiple assignments with different conditions. This feature in Chisel simplify the conversion from \textbf{fsm.assign} to an assignment in Chisel. Pipeline style components need some extra register which indicates the control of pipeline structure.

% \fsm dialect closely resembles actual circuit, so generating Chisel from \fsm{} is quite intuitive. The multiplexer is automatically inserted by Chisel, which makes it possible to multiple assignments with different conditions. In Chisel each signal in \fsm is a variable, while each component is converted to a class. The hybrid representation of \fsm can also be passed out to Chisel programs that the only difference of signal in handshake style is encapsulated in class \texttt{DecoupledIO}. This feature significantly remove handshake signals and improve the readability. 

% All the primitives in \fsm are implemented as chisel templates including arithmetic operations and elastic units. The handshake protocol is encapsulated in class \textbf{DecoupledIO}

\iffalse
\subsection{Partition Pass}
The time graph generated by Static Scheduling Pass can provide useful information for partition pass to explore different design trade-offs. Section \ref{sec:algorithm} describes our partition algorithm used in this pass.

\subsection{Partition Pass}
The time graph generated by Static Scheduling Pass provides 
The time graph structure provide a flexible way to achieve mixed scheduling, in which the time edges represents most of the differences between dynamic and static scheduling. Because dynamic scheduling doesn't need the exact cycle information, the time graph that static scheduling generates can be easily converted to dynamic scheduling by explicitly cleaning up time information. In addition, appropriate time information can indicate the execution model of dynamic modules, and all the operations allocated to the same time state have a high probability to run at the same time. Therefore, these operations can be seen as a whole, which means the partition algorithm will dispute these operations into the same module. Based on this idea, we first use a static scheduling algorithm and then adopt a partition algorithm to convert some edges to dynamic ones. This partition algorithm can be easily applied on \dsir{} representation, which is implemented as transformation on time graph. The main idea of partition algorithm is to minimize the expectation initial interval with the choice of each time edge. The details of partition algorithm is beyond the scope of this paper.

\subsection{Hybrid Scheduling Pass}
After partition is done, we need to re-run static scheduling for each submodule
\subsection{\dsir{} Lowering Pass}
Mixed modules in \dsir{} are split into purely static or dynamic \dsir{} modules in split pass by outlining each connected static scheduling sub-graph as a individual function. Having done some transformations like re-scheduling and eliminating time graph, these static and dynamic modules which are then lowered to the \fsm{} components separately. The static \fsm{} generation pass converts each purely static \dsir{} function into a \fsm{} component, which allocates operations into multiple stages for loop/function pipelining behavior or describes a state transition graph in \fsm{} dialect's syntax otherwise. The conflict graph and live range are analyzed based on the time graph representation, which improves the sharing of resource and register in \fsm{}. To support pipeline structure, the register between register should also be explicitly defined. As for dynamic modules, elastic components are inserted to cope with handshake signal. Each inner sub-module is declared as a \emph{fsm.instance} operation in both static and dynamic modules, and the input of these modules is defined through \emph{fsm.assign}. After doing these transformations, \dsir{} is lowered to \fsm{}.

\subsection{Chisel Generation Pass}
\fsm{} is closer to hardware, which explicitly defines all needed connections, registers and compute units. So it doesn't need many transformations that convert the components in \fsm{} to chisel program. All the primitives in \fsm{} are implemented as a chisel class, which includes arithmetic operations and elastic components. Chisel provides the support for handshake protocol, which is encapsulated in class \textbf{DecoupledIO}. Each \textbf{fsm.primitive} and \textbf{fsm.instance} is converted to a instance. Because the port name is stored in the result of operation, these port can be easily to access like \textbf{adder.result}. \textbf{stateset} region is translated to the finite state machine in Chisel which has the same functionality except the go insertion in the initial state. The multiplexer will be automatically inserted by Chisel, which makes it possible to multiple assignments with different conditions. This feature in Chisel simplify the conversion from \textbf{fsm.assign} to an assignment in Chisel. Pipeline style components need some extra register which indicates the control of pipeline structure.

\fi
\section{Case studies}
We use Hector to build three synthesis approaches: an HLS tool, a systolic array generator, and an accelerator component for sparse linear algebra. These case studies are chosen to show how Hector can be used to build various synthesis tools. %Furthermore, we want to demonstrate that Hector makes it possible to flexibly combine different approaches. The details of these case studies are explained below.

\begin{table*}[htbp!]
    \centering
    
    \caption{Timing and Resource Comparison between Hector and Vitis HLS (Vitis), Dynamatic (DYN)}
    \label{tab:static_dynamic}
    \resizebox{\textwidth}{!}{
    \begin{tabular}{@{} l *{21}{r} @{}}
    \toprule
    \multirow{2}{*}{\textbf{Benchmark}} & \multicolumn{3}{c@{}}{\textbf{Slices}} & \multicolumn{3}{c@{}}{\textbf{LUTs}} & \multicolumn{3}{c@{}}{\textbf{FFs}} & \multicolumn{3}{c@{}}{\textbf{DSPs}} & \multicolumn{3}{c@{}}{\textbf{Cycles}}
    & \multicolumn{3}{c@{}}{\textbf{Period (ns)}} & \multicolumn{3}{c@{}}{\textbf{Time (ms)}} \\
    \cmidrule(lr){2-4} \cmidrule(lr){5-7} \cmidrule(lr){8-10} \cmidrule(lr){11-13} \cmidrule(lr){14-16} \cmidrule(lr){17-19} \cmidrule(lr){20-22}
    
    & Vitis & ours & ratio & Vitis & ours & ratio & Vitis & ours & ratio & Vitis & ours & ratio & Vitis & ours & ratio & Vitis & ours & ratio & Vitis & ours & ratio \\
    \midrule
    \textbf{GEMM} & 486 & 441 & 0.846 & 852 & 890 & 1.045 & 1958	& 1600 & 0.817 & 14 & 13 & 0.929 & 3932.1k & 3752.1k & 0.954 & 5.073 & 4.14 & 0.816 & 19.948 & 15.533 & 0.779 \\
    \textbf{Stencil2D} & 47 & 138 & 2.936 & 94 & 192 & 2.043 & 188 & 370 & 1.968 & 3 & 4 & 1.333 & 320.5k & 312.9k & 0.976 & 4.545 & 3.904 & 0.859 & 1.457 & 1.221 & 0.838 \\
    \textbf{Stencil3D} & 204 & 245 & 1.201 & 454 & 372 & 0.817 & 668 & 890 & 1.332 & 6 & 8 & 1.333 & 102.5k & 103.7k & 1.011 & 5.692 & 4.672 & 0.821 & 0.583 & 0.484	& 0.30 \\
    \textbf{SPMV (CSR)} & 502 & 438 & 0.873 & 881 & 932 & 1.058 & 1934 & 1625 & 0.840 & 14 & 13 & 0.929 & 37.1k & 34.2k & 0.920 & 5.299 & 4.848 & 0.915 & 0.197 & 0.165 & 0.842 \\
    \midrule
    \textbf{Normalized mean} & & & \textbf{1.388} & & & \textbf{1.189} & & & \textbf{1.177} & & & \textbf{0.964} & & & \textbf{0.966} & & & \textbf{0.876} & & & \textbf{0.845} \\
    
    \bottomrule
    
    & DYN & ours & ratio & DYN & ours & ratio & DYN & ours & ratio & DYN & ours & ratio & DYN & ours & ratio & DYN & ours & ratio & DYN & ours & ratio \\
    \midrule
    \textbf{AEloss Pull} & 115 & 99 & 0.860 & 331 & 280 & 0.846 & 265	& 212 & 0.800 & - & - & - & 12.5k & 14.7k & 1.175 & 6.1 & 5.569 & 0.913 & 0.076 & 0.082 & 1.073 \\
    \textbf{AEloss Push} & 365 & 87 & 0.238 & 1118 & 250 & 0.224 & 900 & 199 & 0.221 & - & - & - & 326.1k & 293.7k & 0.901 & 6.2 & 5.5 & 0.887 & 2.022 & 1.616 & 0.799 \\
    \textbf{Stencil2D} & 534 & 408 & 0.764 & 1626 & 1227 & 0.755 & 1379 & 891 & 0.646 & - & - & - & 429.8k & 398.6k & 0.928 & 7.277 & 6.590 & 0.906 & 3.128 & 2.627 & 0.840 \\
    \midrule
    \textbf{Normalized mean} & & & \textbf{0.621} & & & \textbf{0.680} & & & \textbf{0.556} & & & \textbf{-} & & & \textbf{1.001} & & & \textbf{0.901 } & & & \textbf{0.903} \\
    \bottomrule
    \end{tabular}
    }
\end{table*}

\begin{table*}[htbp!]
    \centering
    \resizebox{\textwidth}{!}{
    \begin{tabular}{@{} l *{21}{r} @{}}
    \toprule
    \multirow{2}{*}{\textbf{Benchmark}} & \multicolumn{3}{c@{}}{\textbf{Slices}} & \multicolumn{3}{c@{}}{\textbf{LUTs}} & \multicolumn{3}{c@{}}{\textbf{FFs}} & \multicolumn{3}{c@{}}{\textbf{DSPs}} & \multicolumn{3}{c@{}}{\textbf{Cycles}}
    & \multicolumn{3}{c@{}}{\textbf{Period (ns)}} & \multicolumn{3}{c@{}}{\textbf{Time (ms)}} \\
    \cmidrule(lr){2-4} \cmidrule(lr){5-7} \cmidrule(lr){8-10} \cmidrule(lr){11-13} \cmidrule(lr){14-16} \cmidrule(lr){17-19} \cmidrule(lr){20-22}
    
     & S & D & H & S & D & H & S & D & H  & S & D & H &  S & D & H &S & D & H & S & D & H  \\
    \midrule
    \textbf{AELoss Pull} & 833 & 2001 & 1607 & 1844 & 5216 & 3447 & 3198	& 7193 & 7737 & 16 & 55 & 29 & 15.4k & 9.7k & 9.3k & 5.378 & 5.683 & 5.765 & 0.083 & 0.055 & 0.054 \\
    \textbf{AELoss Push} & 2539 & 2429 & 3058 & 4471 & 7593 & 7339 & 11123 & 8445 & 11543 & 6 & 15 & 12 & 1506.2k & 970.9k & 723.6k & 6.004 & 6.058 & 6.14 & 9.043 & 5.882 & 4.443 \\
    \textbf{Stencil2D} & 138 & 442 & 138 & 192 & 1219 & 192 & 370 & 924 & 370 & 4 & 4 & 4 & 312.9k & 304.9k & 312.9k & 3.904 & 6.4 & 3.904 & 1.221 & 1.951	& 1.221 \\
    \midrule
    \textbf{Normalized mean} & \textbf{1} & \textbf{2.139} & \textbf{1.378} & \textbf{1} & \textbf{3.625} & \textbf{1.503} & \textbf{1} & \textbf{1.835} & \textbf{1.486} & \textbf{1} & \textbf{2.313} & \textbf{1.6044} & \textbf{1} & \textbf{0.748} & \textbf{0.695} & \textbf{1} & \textbf{1.235} & \textbf{1.032} & \textbf{1} & \textbf{0.969} & \textbf{0.713} \\
    
    \bottomrule
    \end{tabular}
    }
    \caption{Timing and Resource Comparison between Hector's static scheduling (S), dynamic scheduling (D) and hybrid scheduling (H)}
    \label{tab:compare}
\end{table*}

\iffalse
\begin{table*}[htbp!]
    \centering
    \resizebox{\textwidth}{!}{
    \begin{tabular}{@{} l *{21}{r} @{}}
    \toprule
    \multirow{2}{*}{\textbf{Benchmark}} & \multicolumn{3}{c@{}}{\textbf{Slices}} & \multicolumn{3}{c@{}}{\textbf{LUTs}} & \multicolumn{3}{c@{}}{\textbf{FFs}} & \multicolumn{3}{c@{}}{\textbf{DSPs}} & \multicolumn{3}{c@{}}{\textbf{Cycles}}
    & \multicolumn{3}{c@{}}{\textbf{Period (ns)}} & \multicolumn{3}{c@{}}{\textbf{Time (ms)}} \\
    \cmidrule(lr){2-4} \cmidrule(lr){5-7} \cmidrule(lr){8-10} \cmidrule(lr){11-13} \cmidrule(lr){14-16} \cmidrule(lr){17-19} \cmidrule(lr){20-22}
    
    & Dynamatic & ours & ratio & Dynamatic & ours & ratio & Dynamatic & ours & ratio & Dynamatic & ours & ratio & Dynamatic & ours & ratio & Dynamatic & ours & ratio & Dynamatic & ours & ratio \\
    \midrule
    \textbf{AEloss Pull} & 115 & 99 & 0.860 & 331 & 280 & 0.846 & 265	& 212 & 0.800 & - & - & - & 12.5k & 14.7k & 1.175 & 6.1 & 5.569 & 0.913 & 0.076 & 0.082 & 1.073 \\
    \textbf{AEloss Push} & 365 & 87 & 0.238 & 1118 & 250 & 0.224 & 900 & 199 & 0.8 & - & - & - & 326.1k & 293.7k & 0.901 & 6.2 & 5.5 & 0.887 & 2.022 & 1.616 & 0.799 \\
    \textbf{Stencil2D} & 534 & 408 & 0.764 & 1626 & 1227 & 0.755 & 1379 & 891 & 0.646 & - & - & - & 429.8k & 398.6k & 0.928 & 7.277 & 6.590 & 0.906 & 3.128 & 2.627 & 0.840 \\
    \midrule
    \textbf{Normalized mean} & \textbf{1$\times$} & \textbf{2.313$\times$} & \textbf{1.604$\times$} & & & \textbf{0.680$\times$} & & & \textbf{0.556$\times$} & & & \textbf{-} & & & \textbf{1.001$\times$} & & & \textbf{0.901 $\times$} & & & \textbf{0.903$\times$} \\
    \bottomrule
    \end{tabular}
    }
    \caption{Timing and Resource Comparison between Hector and Dynamatic}
    \label{tab:dynamic}
\end{table*}
\fi

\subsection{High-level Synthesis}
\label{sec:hls}
% In this section, we show the convenience and extensibility of the proposed flow as well as the profits brought by hybrid scheduling method. 

% To prove that Hector can generate comparable hardware, we first evaluate Hector on several benchmarks compared with static and dynamic HLS tools separately. Then We compare the quality of three kinds of scheduling strategies, fully static scheduling, fully dynamic scheduling, and a hybrid scheduling method in Hector, which demonstrates Hector provides flexibility to explore design trade-offs.



Recently, \cite{cheng_dss} demonstrates the effectiveness of hybrid scheduling. Their approach manually partitions a program into static and dynamic regions and then synthesizes them using either static or dynamic scheduling HLS tools, respectively. Then, they glue different regions together at the register-transfer level. Their later work~\cite{find_islands} tries to avoid manual selection by implementing a simple partition algorithm in LLVM IR. However, LLVM IR is mainly a software compilation IR and lacks hardware semantics. Therefore, one central question is \textit{what are the intermediate abstractions suitable for both static and dynamic schedulings}?

In this case study, we build an HLS tool based on Hector to demonstrate the generality of the proposed IR. MLIR provides several built-in Dialects, such as Affine, SCF, and Standard~\cite{mlir}. SCF Dialect describes static control flow in a higher-level abstraction which is selected as the input of HLS. To build a complete HLS flow, we additionally implement an SDC-based module scheduling~\cite{sdc-modulo-13} that converts SCF Dialect to \tor. The scheduling result can be easily represented by \tor because of the concise structure of the time graph. Hybrid scheduling is also supported as a hybrid time graph which contains multiple manners in the same graph. We use a naive partition strategy on \tor that converts the manners of some edges from static to dynamic. This demonstrates that Hector can flexibly support different scheduling strategies.

\textbf{Benchmarks}. We evaluate using both regular and irregular applications. The regular kernels are selected from Machsuite\cite{machsuite} including GEMM, Stencil2D, Stencil3D, and Spmv(CSR). These four kernels do not have any branch, and thus are well suited for static scheduling. The irregular kernels are AELoss Pull and AELoss Push from the loss function of~\cite{aeloss}. For these two kernels, a loop-carried dependence resides inside the innermost loop. A considerable amount of computation is guarded by an unpredictable condition, which is seldom true.

\textbf{Methodology}. Polygeist\cite{polygeist} is used as Hector's front-end to convert C programs into SCF dialect in MLIR, then Hector generates an RTL file in Verilog targeting FPGA part xc7z020clg484. We obtain cycle numbers using RTL level simulation. Timing and resource results are obtained from post implementation report from Vivado 2021.1.

In the following, we first compare Hector with the state-of-the-art HLS frameworks. Specifically, we compare static scheduling in Hector with Vitis HLS~\cite{VivadoHls} and dynamic scheduling with Dynamatic~\cite{Dynamatic_Josipovi} \footnote{Dynamatic only provides one buffer insertion strategy based on a Mix ILP solver, which is time-consuming and unscalable. Buffer insertion is disabled in both tools for a fair comparison.}. For Vitis HLS, we apply pipeline directives on the inner-most loop. For comparison with Dynamatic, we convert the floating-point operations for benchmarks AELoss Pull and AELoss Push to their integer equivalent because the open-source Dynamatic with floating-point computation fails to run. Then, we evaluate the effect of our hybrid scheduling in Hector by comparing it with static and dynamic scheduling. In Hector, we can switch between static, dynamic, and hybrid scheduling by easily configuring the scheduling modes provided as an attribute of SCF functions.

\textbf{Comparison with state-of-the-art HLS.} \autoref{tab:static_dynamic} demonstrates that Hector can achieve comparable performance with Vitis HLS and Dynamatic for the tested kernels. The number of DSP in Vitis HLS is different because we use a different IP configuration. For Stencil2D and Stencil3D benchmarks, we attribute the large consumption of resources to our resource binding algorithm compared with Vitis HLS. Hector generates smaller hardware than Dynamatic because of Chisel's internal optimizations, such as width analysis. For the AELossPush benchmark, the resource usage of Dynamatic is substantially higher because LLVM IR, the input language of Dynamatic, has too many basic blocks, which have a negative impact on Dynamatic's basic blocks based scheduling algorithm.

\textbf{Hybrid scheduling.} For the AELossPull benchmark, the major bottleneck is caused by the loop-carried dependence guarded by a condition. The loop-carried dependence causes a conservative II in static scheduling, while dynamic scheduling yields better throughput. Our partition algorithm extracts out the operations except the ones associated with the loop-carried dependence into a static submodule. The submodule contains four floating-point multiplications and four floating-point additions, which can be implemented by only two multipliers and two adders in static scheduling. However, these operations can not share any resources in dynamic scheduling, leading to massive DSP consumption. The hybrid scheduling can retain the performance improvement of dynamic scheduling while reducing resource usage by enabling resource sharing inside static submodules. Static (dynamic) scheduling is effective for pure regular (irregular) applications, while hybrid scheduling is a promising solution for complex applications with a mixture of regular and irregular behaviors. Overall, the hybrid scheduling improves the performance by 29$\%$ compared to static on average. 

AELoss Push contains a two-level loop nest with conditions. The comparison results between the three scheduling modes are similar to AELoss Pull except for the higher FF usage in static and hybrid scheduling. This is mainly because of our register allocation strategy. The complicated loop nest and conditions exacerbate the buffer insertion algorithm's effect, while the regular behavior mitigates this situation in static scheduling. As a result, hybrid scheduling preserves the irregular behavior and achieves a better throughput with the help of static submodules. As for the Stencil2D benchmark, the partitioning algorithm chooses the whole program as a static module because the computation has no branches and is very regular. In that case, dynamic scheduling can only get comparable performance while consuming a large number of resources.

The comparison result against existing HLS tools demonstrates that Hector can fit well in HLS no matter with static, dynamic, or hybrid manners.~\cite{find_islands} implements a hybrid scheduling strategy by partitioning LLVM IR into multiple functions and synthesizing them with separate HLS tools. In contrast, Hector allows users to customize HLS algorithms (such as the partitioning algorithm) implemented as MLIR passes, which facilitates easy extension and allows more optimization opportunities.

% \hl{the result looks good, but we do not discuss how the IRs are used for different scheduling.}
\subsection{Systolic Array Generator}

%Data reuse occurs within the PE array, avoiding broadcasting and resulting in low fan-ins. Systolic array achieves state-of-the-art performance in many domains, like machine learning~\cite{TPU}, due to its high frequency. \autoref{fig:systolic}(a) depicts an example of a systolic array accelerator for matrix multiplication, in which two input matrices are reused through the interconnection between PE array. Generating hardware written in RTL programs is difficult due to the low level of abstraction. However, generating a systolic array at a high-level abstraction is still complicated. Because the HLS procedure often ignores the architecture information, inserting suitable directives is required for an HLS implementation, which is tough even for an expert~\cite{autoSA}. 

Systolic array is a class of accelerators where inputs or results pass through the processor elements (PEs)~\cite{polySA, autoSA, tensorlib, ems}. Despite the fact that systolic array offers good performance, it's challenging to conveniently generate a customized accelerator. In this case study, we use Hector to build a systolic array generator. This generator can generate PE arrays of arbitrary size given a PE's configuration, which can be written in either \tor or \hec. In addition, \texttt{SCF} Dialect is also supported by reusing the previous HLS process (\autoref{sec:hls}).

\begin{figure}[htbp!]
  \centering
  \includegraphics[width=\linewidth]{figure/code.pdf}
  \caption{A 2$\times$n by n$\times$2 matrix multiplication systolic array described in \hec. The interconnection between PEs is described as \texttt{hec.assign}.}
  \label{fig:systolic}
  \vspace{-9pt}
\end{figure}

\textbf{Implementation.} We use an elastic implementation that allows each PE to stall if the data does not arrive or the subsequent PE is not ready to receive new data. Supporting dynamic behavior is critical, particularly when bandwidth is limited. When sufficient bandwidth is ensured, an elastic systolic array performs similarly to the static systolic array with the exception of a few more signals. This dynamic behavior is implemented as an elastic interconnection between PE array. \autoref{fig:systolic} shows a 2$\times$2 systolic array represented in \hec. The outer module is a dynamic style component that enables the implementation to be elastic. All PEs are instantiated from the corresponding component, and data movement between PEs is explicitly stated as \verb|hec.assign|.

A single PE can be configured at multiple levels. For the PE with simple behavior, \texttt{SCF} Dialect can be automatically lowered using the approach described in \autoref{sec:hls}. \hec can be used in applications with no resource limitations, and \tor is capable of handling complex designs. This once again demonstrates the advantages of multi-level IR. Pipelined PE can be easily expressed at these levels due to the explicit pipeline in both levels.

\begin{figure}[htbp!]
  \centering
  \includegraphics[width=\linewidth]{figure/data_systolic.pdf}
  \caption{Cycle count comparison between different implementations of matrix multiplication.}
  \label{fig:data_systolic}
  \vspace{-9pt}
\end{figure}


\textbf{Comparison against Calyx.} Calyx~\cite{} also provides a static systolic array generator for demonstration purposes, but it lacks the expressivity required for hardware generation. More than 3000 lines of code are required to describe an 8×8 systolic array. The fundamental reason is that the control logic is described in a flattened representation that explicitly represents all data moves at each cycle. Another issue in Calyx is that it lacks a pipeline primitive, prohibiting all PEs in generated hardware from being pipelined. In real-world applications, interleaving multiple calculations is commonly used to improve resource utilization ratio, but the lack of pipeline primitives makes this optimization impossible. 

\textbf{Evaluation.} We evaluate matrix-multiplication kernels ranging from 4×4 to 32×32. The PE configuration of the systolic array is written in \texttt{SCF}, which is automatically lowered to \hec. Aside from the systolic arrays generated by Hector and Calyx, we implement a basic program in Vitis HLS that uses pipeline pragma in the inner-most loop. We optimize the HLS implementation by unrolling the outer two loops and partitioning all the matrices in the proper dimension to make a fair comparison. We also evaluate Tensorlib~\cite{tensorlib}, a state-of-the-art systolic array generator, with the same dataflow configuration. The experiment targets Zynq UltraScale+ XCZU3EG FPGA at a 7ns clock period, which is the same as Calyx's setup. We compare the cycle counts of the designs (\autoref{fig:data_systolic}). Hector supports pipeline PE, which further optimizes the performance by interleaving multiple computations. On average, the systolic array generated by Hector improves the performance by 5.6× improvement compared to optimized HLS implementation. In comparison to the HLS implementation, Calyx implementation with sequential PEs can only achieve 79\% performance, demonstrating the importance of pipeline semantics. Hector can achieve the same performance with Tensorlib due to structural description and pipeline support.

The experiment results demonstrate Hector's versatility in applying domain-specific approaches, which is difficult to express in HLS tools. With architectural information, which is allowed in Hector, it's straightforward to outperform a general-purpose HLS tool. The multi-level paradigm also provides the ability to design hardware at various levels.

\subsection{Shuffle Unit}
\label{sec:spmv}
%Some studies concentrate on architectural design and implement various optimizations to improve throughput. However, due to the usage of sophisticated cache and crossbar structures, only a few works are compatible with FPGA.

%Sparse linear algebra is now widely used in various domains such as scientific computation and machine learning.  Implementing an accelerator for sparse linear algebra at the RTL level is extremely challenging~\cite{spaghetti}. Here, we select sparse matrix-vector multiplication (SPMV) as a case study to exhibit the extensibility of Hector. Due to the characteristics of SPMV like memory bound and high sparsity, a streaming accelerator fits well in the sparse computation. The dynamic behavior leads to complex control of stalling, which increases the difficulty of designing at the RTL level. An HLS implementation for SPMV has been proposed in some works. Alberto Parravicini et al.~\cite{topK} store a redundancy vector for each PE, which is unscalable for large-scale due to the limited on-board memory. To deal with the limited on-board memory, it's necessary to use a shared buffer between different PEs. 

Shuffle unit is the key hardware component that solves bank conflicts between multiple PEs, which has been used in sparse matrix-vector multiplication (SPMV) design. Each pair of streams may access an arbitrary bank determined by the column index, resulting in traffic conflict. However, it is challenging to implement the shuffle unit with optimized throughput in HLS tools~\cite{HiSparse}. In static scheduling, the conservative assumption of traffic pattern results in a big II, and dynamic scheduling is not suited to generate efficient control logic. Therefore, it's better to implement the shuffle unit at the RTL level. %Designing important modules at the RTL level and quickly generating secondary modules is a better solution for SPMV.

In Hector, the combination of RTL and other synthesis approaches is simple. Given a dummy representation with corresponding ports in \hec and \tor, other components can easily instantiate it ignoring the inner structure. All of the transformations only work on those real components, and the dummy component is eventually replaced with the Chisel program. RTL implementation can be easily integrated with other synthesis approaches using the Chisel implementation and dummy representations.

% \definecolor{bg}{rgb}{0.95,0.95,0.95}
% \begin{minted}[bgcolor=bg,fontsize=\footnotesize]{c}
% tor.module Shuffle(in0, id0, in1, id1)->(out0, out1) {
% }{"dynamic", "dummy"} //tor dummy representation
% hec.component Shuffle(in0, id0, in1, id1, out0, out1) {
% }{"dynamic", "dummy"} //hec dummy representation
% class Shuffle extends MultiIOModule {
% } //Hardware logic of Shuffle Unit
% \end{minted}

\begin{table}[htbp!]
\footnotesize
\renewcommand\arraystretch{1.5}
% \setlength{\tabcolstep}{20mm}
\begin{tabular}{lcc}
\hline
\multicolumn{1}{c}{} & Line of Codes & Initial Interval \\ \hline
Hector (RTL)         & 35            & 1                \\ \hline
Vitis HLS            & 111           & 1                \\ \hline
\end{tabular}
\caption{The comparison between Vitis HLS and Hector.}
\label{tab:shuffle}
\vspace{-9pt}
\end{table}

%The throughput of a streaming accelerator is governed by the weakest component, hence a poor component degrades overall performance. 
%The shuffle unit, as a critical hardware module in SPMV, is compared against HLS and RTL implementation separately. 
We use Hector to construct an RTL design and simplify a pipeline kernel~\cite{HiSparse} written in Vitis HLS. The RTL implementation takes only 35 lines to implement the control logic, while the optimal HLS implements uses 111 lines. As previously stated, the conservative assumption in HLS results in a large II. Due to explicit control logic in HLS implementation, it get the optimal II which matches the RTL implementation. The experiment demonstrates that HLS is not suitable for control logic. Therefore, combining HLS and RTL design is a superior option for achieving good performance while maintaining enough productivity.

% \hl{let us include a table for result}

% \begin{figure}[htbp!]
%   \centering
%   \begin{subfigure}[b]{0.45\linewidth}
%   \centering
%   \includegraphics[width=\linewidth]{figure/cycle_spmv.pdf}
%   \caption{Cycle count}
%   \end{subfigure}
%   \hfill
%   \begin{subfigure}[b]{0.45\linewidth}
%   \centering
%   \includegraphics[width=\linewidth]{figure/lut_spmv.pdf}
%   \caption{LUT usage}
%   \end{subfigure}

%   \caption{Performance and resource comparison between HLS and RTL implementation in Hector.}
%   \label{fig:spmv}
% \end{figure}
% \subsection{Discussion}
% These case studies show a variety of synthesis techniques, including different HLS scheduling strategies, domain-specific approaches, and a hybrid HLS/RTL design. All of these methods can be well supported in Hector with minimal engineering effort. Hector also allows for the flexible combination of multiple methodologies, which is difficult to achieve at the RTL level.

\section{Related Works}

In recent years, hardware intermediate representations design has become a popular topic. An appropriate IR can give sufficient expressivity and avoid redundant engineering efforts. Several existing hardware IRs are introduced next.

\paragraph{Domain-specific languages} DSLs provide a higher level of abstraction that is natural to the domain, assisting in the generation of hardware with expert knowledge. Raghu Prabhakar et al.~\cite{parallel_pattern} adopt parallel patterns like \emph{map} and \emph{fold} to express high-level computation. Aetherling~\cite{aetherling} aims at generating streaming accelerators by applying transformations to the proposed data-parallel IR. Nithin George et al.~\cite{system_dsl} generate hardware systems from applications written in a machine-learning DSL. HeteroCL~\cite{heterocl} provides an abstraction that decouples algorithm description and architecture specification, and develops a compilation flow to heterogeneous computing platforms. Apart from domain-specific optimizations, extra compilation to the RTL program is required, which makes it difficult for researchers to generate hardware quickly.

% \hl{can we merge the IRs into one paragraph?}

% \paragraph{IR for HDL} FIRRTL~\cite{firrtl} that obeys AST format is the intermediate representation in Chisel. It is a multi-level IR in which high-level abstraction is lowered by simplifying and reducing some high-level semantics like some object-oriented features. LLHD\cite{LLHD_Schuiki} is a three-level IR that includes behavioral, structural, and netlist levels of hardware synthesis. Different level in LLHD aims at different applications including simulation, verification, and logic synthesis. Hector uses higher-level abstraction to combine different synthesis methods and avoid redundant implementations, whereas these IRs are still at a low-level abstraction with no high-level information.

% \paragraph{IR for HLS} Many HLS tools~\cite{VivadoHls, legup, Dynamatic_Josipovi} still adopt LLVM IR as their internal representation, which is a pure software IR that doesn't contain any hardware information. $\mu$IR~\cite{uIR_Sharifian} uses dynamic scheduling based on task-level parallelism and briefly presents a task-level representation. Wu et al.~\cite{hierarchy_cdfg} propose a hierarchical CDFG representation for hardware and software co-design. SynASM~\cite{synASM} proposes a timed CDFG that is extended with some timing information such as sequential and parallel, and transforms it directly to VHDL using its own HLS process. AHIR~\cite{ahir} proposes a low-level abstraction that decouples the datapath and control path. All of these IRs only focus on one single abstraction level which loses the opportunity of interfering with HLS at different levels. Furthermore, these representations aren't able to represent hardware with both static and dynamic behavior.

% \paragraph{Hardware infrastructure} Calyx\cite{Calyx_Rachit} provides software-like control flow primitives such as seq, par, and while, that are then lowered to the finite state machine (FSM). It's not suitable for high-level synthesis or hardware with complex pipelines since it lacks pipeline and task-level semantics, resulting in poor expressivity. Moreover, Calyx only supports FSM-controlled hardware, so hardware with dynamic behavior like streaming accelerator is not supported. The goal of CIRCT~\cite{circt} is to construct a reusable and modular infrastructure for the entire hardware generation including high-level synthesis and logic synthesis. This project is still in progress and absorbs existing IR designs like Calyx and LLHD. In contrast, Hector concentrates on a higher abstraction than RTL, which can be regarded as the front-end of hardware generation.


\paragraph{IR for hardware} FIRRTL~\cite{firrtl} that obeys AST format is the intermediate representation in Chisel. LLHD~\cite{LLHD_Schuiki} is a three-level IR that aims at different applications including simulation, verification, and logic synthesis. $\mu$IR~\cite{uIR_Sharifian} uses dynamic scheduling based on task-level parallelism and briefly presents a task-level representation. Wu et al. and synASM~\cite{hierarchy_cdfg, synASM} propose a CDFG represention for hardware and software. AHIR~\cite{ahir} proposes a low-level abstraction that decouples the datapath and control path. Calyx\cite{Calyx_Rachit} provides software-like control flow primitives such as seq, par, and while to describe hardware. The goal of CIRCT~\cite{circt} is to construct a reusable and modular infrastructure for the entire hardware generation including high-level synthesis and logic synthesis. This project is still in progress and absorbs existing IR designs like Calyx and LLHD.
\section{Conclusion}
% In this paper, we propose Hector, a holistic HLS framework with versatile scheduling support for FPGAs. Hector built on a multi-level IR paradigm  supports various HLS scheduling strategires including static, dynamic, and hybrid scheduling. We also develop a heuristic partitioning algorithm inside the IR system to automate the hybrid methodology. Experiments show that Hector generates comparable hardware designs with existing HLS tools in terms of performance and resource.  Moreover, the open-source framework provides flexibility to customize synthesis algorithms and allows users to explore advanced HLS scheduling techniques.

In this paper, we propose Hector, a two-level IR providing a unified description for different synthesis methodologies. Through a series of transformations and optimizations based on the MLIR infrastructure, Hector's IRs are finally converted to synthesizable RTL programs. We demonstrate the expressivity and effectiveness of our design by implementing three synthesis approaches. The experiment results show that Hector can generate comparable hardware designs with existing HLS tools in terms of performance and resource utilization. With the combination of different methodologies, it's simple for Hector to outperform HLS tools. Moreover, the open-source framework provides flexibility to customize synthesis approaches and allows users to explore advanced techniques.

\section{Acknowledgments}
This work is supported in part by National Natural Science Foundation of China (NSFC) under grant No.62090021.

%%
%% The next two lines define the bibliography style to be used, and
%% the bibliography file.
\bibliographystyle{ACM-Reference-Format}
% \bibliographystyle{plain}
\bibliography{bibfile}

\end{document}
\endinput
%%
%% End of file `sample-sigconf.tex'.
